% =====================================================================
%  theory/energy_flux_derivation.tex
%  Phase 1, Blueprint V5.0, Task F
%
%  Lord-Shulman energy identity, normal total-field flux, bulk and interface
%  dissipation, positivity, and the power balance that must be established
%  BEFORE R + T + A = 1 may be asserted.
%
%  SOURCE-ONLY: no TeX engine is installed. Verified as text; the flux
%  functional is checked by two independent numerical routes in
%  code/foundation_checks.py (check D1) and its conservation along z in
%  check E1.
% =====================================================================
\documentclass[11pt]{article}
\usepackage{amsmath,amssymb,amsthm,bm}
\usepackage{geometry}
\usepackage{booktabs}
\geometry{margin=1in}
\newcommand{\s}{\sigma}
\newcommand{\eps}{\varepsilon}
\newcommand{\bth}{\vartheta}
\newcommand{\D}{\bm{D}}
\newcommand{\E}{\bm{E}}
\newcommand{\bu}{\bm{u}}
\newcommand{\bv}{\bm{v}}
\newcommand{\bq}{\bm{q}}
\newcommand{\Lam}{\Lambda}
\newcommand{\bbar}[1]{\bar{#1}}
\newcommand{\jj}{[\![}
\newcommand{\kk}{]\!]}
\newcommand{\jmp}[1]{[\![#1]\!]}
\theoremstyle{plain}
\newtheorem{proposition}{Proposition}
\newtheorem{remark}{Remark}

\title{Energy flux, bulk and interface dissipation, and the power balance\\
\large Theory note 4 of 4 --- Phase 1, Blueprint V5.0}
\author{Derived for the Piezothermoelastic Anti-Reflection programme}
\date{4 November 2026}

\begin{document}
\maketitle

\begin{abstract}
Blueprint V5.0 makes the energy account a compulsory gate: the identity
$R+T+A=1$ is explicitly \emph{not} a premise, and gates G4 and G5 block any
energy-ratio claim until the flux and the dissipations are derived. This note
supplies that derivation for the Lord--Shulman piezothermoelastic system, and
\textbf{corrects the flux functional stated in Phase 1}. The corrected
cycle-averaged normal flux is
$\langle J_z\rangle = -\tfrac12\operatorname{Re}[\s_{xz}v_x^* + \s_{zz}v_z^* +
i\omega D_z\varphi^*] + \tfrac12\operatorname{Re}[q_z\Theta^*]/T_0$; a bulk
dissipation that
is positive semidefinite under the stability conditions; interface dissipations
matching the contact laws of the companion note; the composite balance
\eqref{eq:master}; and a clear statement of which terms must be evaluated where
before any partition may be quoted. \textbf{No reflection, transmission or
absorption coefficient is computed in this note, and no anti-reflection claim is
made.}
\end{abstract}

\section{Starting point and conventions}
\label{sec:start}

The bulk system is that of \texttt{theory/governing\_equations.tex}. Only the
relations actually differentiated below are repeated here, so that this note is
self-contained. With $\Lam \equiv 1-i\omega\tau_q$,
\begin{equation}
  \s_{ij} = C_{ijkl}\eps_{kl} - e_{kij}E_k - \beta_{ij}\bth, \qquad
  D_i = e_{ikl}\eps_{kl} + \eps_{ik}E_k + p_i\bth, \qquad
  E_i = -\varphi_{,i},
  \label{eq:constitutive}
\end{equation}
\begin{equation}
  \Lam\,\bq = -\kappa\nabla\bth, \qquad
  \Lam\Big[\rho c_E\dot\bth + T_0\big(\beta_{ij}\dot\eps_{ij} + p_i\dot E_i\big)\Big]
  = \kappa_{ij}\bth_{,ij},
  \label{eq:fluxlaw}
\end{equation}
\begin{equation}
  \rho>0,\quad c_E>0,\quad T_0>0,\quad \kappa_{ij}\ \text{positive definite},\quad
  \tau_q\ge 0,\quad C_{ijkl}\ \text{positive definite},
  \label{eq:thermal-stability}
\end{equation}
together with the momentum balance $\s_{ij,j}=\rho\ddot u_i$ and Gauss'
law $D_{i,i}=0$. \textsc{Derived} / \textsc{Assumed} as labelled in the
governing-equations note. Time dependence $e^{-i\omega t}$; complex amplitudes are
phasors, so physical fields are $\operatorname{Re}(\cdot\,e^{-i\omega t})$ and the
cycle average of a product of two phasors $A,B$ is
$\tfrac12\operatorname{Re}(A B^*)$.

\section{The balance law, and why the thermal term is not $\bm{q}$ alone}
\label{sec:balance-law}

\textbf{Phase 1B.} The Phase-1 version of this section wrote the cycle-averaged
normal flux as
\begin{equation}
  \langle P_z\rangle \overset{\text{Phase 1}}{=}
  -\tfrac12\operatorname{Re}\big[\s_{xz}v_x^*+\s_{zz}v_z^*
  +i\omega D_z\varphi^*\big] + \tfrac12\operatorname{Re}\big[q_z\big],
  \label{eq:phase1-wrong}
\end{equation}
where $q_z$ is the first-order harmonic heat-flux phasor. \textbf{This is not a
valid cycle-averaged energy flux, and the term is withdrawn.} There are two
independent reasons, and both are worth stating because either alone is
sufficient to reject it.

\begin{enumerate}
  \item \textbf{It is linear, not quadratic.} A cycle average
        $\tfrac12\operatorname{Re}(AB^*)$ is a bilinear form in two first-order
        fields. The term $\tfrac12\operatorname{Re}[q_z]$ is linear in one
        field. It cannot be produced by averaging any quadratic flux over a
        period: the average of a monochromatic first-order quantity is
        \emph{zero}, and the average of a quadratic form is nonzero precisely
        because it is a cross term between $q_z$ and something else.
  \item \textbf{It omits the temperature.} The heat flux enters the energy
        budget multiplied by a temperature. Without a second field the term has
        no way to know the state of the medium: it would be invariant under
        $\bth\to\bth+\text{const}$, which is not a symmetry of a conducting body.
\end{enumerate}

\textbf{The correct object.} Work with the \emph{total} energy current, not with
the mechanical Poynting vector plus a bare heat flux. The mechanical and
electrical terms already have the right form; the thermal field contributes
through the convected internal energy plus conduction. Expand the internal energy
about $T_0$ as $u = \rho c_E\bth$ and let $\bth$ be the first-order temperature
perturbation; the absolute temperature is $T=T_0+\bth$. The total energy current
in the $z$ direction is
\begin{equation}
  J_z = \underbrace{-\s_{xz}v_x-\s_{zz}v_z-D_z\dot\varphi}_{\text{mechanical + electrical}}
      + \underbrace{q_z\,\frac{\bth}{T_0}}_{\text{thermal}},
  \label{eq:total-current}
\end{equation}
\textsc{Derived} as follows. The first-order energy balance is
\begin{equation}
  \partial_t\Big[\tfrac12\rho\dot u_i\dot u_i+\tfrac12\s_{ij}\eps_{ij}
    +\tfrac12\E\cdot\D+\tfrac12\tfrac{\rho c_E}{T_0}\bth^2\Big]
  +\nabla\cdot\bm{J} = -\mathcal{D},
  \label{eq:local-balance}
\end{equation}
and matching $\partial_t\big(\tfrac12\rho c_E\bth^2/T_0\big)
= (\rho c_E/T_0)\bth\dot\bth$ against the entropy balance
$T_0\partial_t(\rho\eta)=-\nabla\cdot\bq$ identifies the thermal part of $\bm{J}$
as $\bq\,\bth/T_0$ to first order. The factor $1/T_0$ is forced by the
$1/T_0$ in the stored thermal energy and by $\rho\eta$ carrying $\rho c_E\bth/T_0$.

\subsection{The corrected cycle-averaged normal flux}

\begin{proposition}[Normal total-field flux]
\label{prop:flux}
For a single partial mode with phasor amplitudes
$\bm{a}=(U,W,\Phi,\Theta)^{\mathsf T}$, $e^{-i\omega t}$, and
$\Lam=1-i\omega\tau_q$,
\begin{equation}
  \boxed{\ \langle J_z\rangle
  = -\tfrac12\operatorname{Re}\big[\s_{xz}v_x^{*}+\s_{zz}v_z^{*}
    +i\omega D_z\varphi^{*}\big]
  +\tfrac12\operatorname{Re}\big[q_z\,\Theta^{*}\big]\big/T_0\ }
  \label{eq:Jz}
\end{equation}
with $v_x^{*}=i\omega U^{*}$, $v_z^{*}=i\omega W^{*}$ and
$q_z=-\kappa_{33}(i\xi\Theta)/\Lam$.
\end{proposition}
\begin{proof}
Insert the phasors into \eqref{eq:total-current} and take the cycle average
$\langle AB\rangle=\tfrac12\operatorname{Re}(AB^*)$. The mechanical and
electrical terms reproduce the first bracket. The thermal term is
$\langle q_z\bth\rangle=\tfrac12\operatorname{Re}(q_z\Theta^*)$.
\end{proof}

\textsc{Tested}: three independent statements, all in
\texttt{code/foundation\_checks.py}.

\begin{center}
\begin{tabular}{@{}lll@{}}
\toprule
Test & Reference & Measured \\
\midrule
Uncoupled elastic limit & $\tfrac12\rho c_L\omega^2\lvert W\rvert^2$ (analytic)
  & $1.95\times10^{-16}$ \\
Closed form vs.\ direct cycle average of the \emph{real} fields
  & independent route & $1.2\times10^{-40}$ (Al6061) \\
Antisymmetry $\langle J_z\rangle(\xi)=-\langle J_z\rangle(-\xi)$
  & exact symmetry & $0$ (exact) \\
Local balance $-2\operatorname{Im}\xi\,\langle J_z\rangle=\langle\mathcal{D}\rangle$
  & Second Law & $7.1\times10^{-17}$ (see \S\ref{sec:balance-verify}) \\
\bottomrule
\end{tabular}
\end{center}

\subsection{Size of the Phase-1 error}
\label{sec:error-size}

The two forms are not close. For the thermal branch of PZT-5H at 1\,MHz the
thermal contribution is
\begin{center}
\begin{tabular}{@{}lrr@{}}
\toprule
& Phase 1: $\tfrac12\operatorname{Re}[q_z]$ & Corrected: $\tfrac12\operatorname{Re}[q_z\Theta^*]/T_0$ \\
\midrule
value ($\mathrm{W\,m^{-2}}$) & $1.99\times10^{31}$ & $-1.30\times10^{3}$ \\
\bottomrule
\end{tabular}
\end{center}
a discrepancy of $34$ orders of magnitude, because the uncorrected term scales
with the flux alone while the corrected one is suppressed by the thermal
relaxation and by $1/T_0$. On the propagating branches the uncorrected term is
$\mathcal{O}(10^{18})$ against mechanical terms of $\mathcal{O}(10^{20})$; it does
not merely shift the answer, it corrupts it at the percent level and, for the
thermal branch, destroys it entirely. \textbf{Any result computed with the Phase-1
formula is void.} None was reported.

\section{Verification of the balance law}
\label{sec:balance-verify}

The flux formula \eqref{eq:Jz} is verified against the Second Law rather than
against itself. For a mode,
$\langle J_z\rangle(z)=\langle J_z\rangle(0)e^{-2\operatorname{Im}(\xi)z}$, so
$\mathrm{d}\langle J_z\rangle/\mathrm{d}z=-2\operatorname{Im}(\xi)\langle J_z\rangle$.
The local balance \eqref{eq:local-balance} therefore requires
\begin{equation}
  -2\operatorname{Im}(\xi)\,\langle J_z\rangle
  \;=\;\langle\mathcal{D}\rangle ,
  \label{eq:balance-test}
\end{equation}
and the dissipation is written directly from the entropy production
$\sigma_s=\bq\cdot\nabla(1/T)=-\bq\cdot\nabla\bth/T^2$, so that
\begin{equation}
  \langle\mathcal{D}\rangle
  = -\frac{1}{T_0}\Big[\tfrac12\operatorname{Re}\big(q_z(\partial_z\bth)^*\big)
    +\tfrac12\operatorname{Re}\big(q_x(\partial_x\bth)^*\big)\Big],
  \label{eq:D-both}
\end{equation}
with $\partial_z\bth=i\xi\Theta$ and $\partial_x\bth=ik\Theta$. \textbf{Equation
\eqref{eq:D-both} shares no amplitude with the mechanical or electrical parts of
\eqref{eq:Jz}}, so \eqref{eq:balance-test} is a genuine test of the functional.

\begin{proposition}[Causality sign of the thermal term]
\label{prop:thermal-sign}
The sign of the thermal term in \eqref{eq:Jz} is fixed by requiring
$\langle\mathcal{D}\rangle\ge0$. With $q_z=\Lam^{-1}(-\kappa_{33}i\xi\Theta)$ and
$\Lam=1-i\omega\tau_q$,
\begin{equation}
  -\tfrac12\operatorname{Re}\big[q_z(i\xi\Theta)^*\big]
  =\tfrac12\frac{\kappa_{33}\lvert\xi\Theta\rvert^{2}}{1+\omega^2\tau_q^{2}}
  \;\ge\;0 ,
  \label{eq:D-analytic}
\end{equation}
independently of the sign convention for $\Lam$, because only the real part of
$\Lam^{-1}$ enters. The opposite choice of sign in \eqref{eq:Jz} would give
$\langle\mathcal{D}\rangle\le0$ and an exponentially \emph{growing} flux along
$-z$, violating the Second Law.
\end{proposition}
\textsc{Tested}: \eqref{eq:D-analytic} is reproduced to $1.19\times10^{-41}$
relative.

\begin{remark}[The transverse term is not optional]
\label{rem:transverse}
Dropping the $q_x\partial_x\bth$ term of \eqref{eq:D-both} leaves a relative
residual in \eqref{eq:balance-test} of exactly $(\kappa/\xi)^2$
$=1.2\times10^{-7}$ at $30^\circ$ for PZT-5H. This is a geometric factor, not a
defect in the thermal term, and it was the reason a first attempt at this check
appeared to fail by a factor consistent with a missing term. Both directions of
conduction must be carried. \textsc{Tested}: including it reduces the residual
from $1.86\times10^{-5}$ to $7.1\times10^{-17}$.
\end{remark}

\begin{remark}[Status of this verification]
\label{rem:balance-status}
The identity \eqref{eq:balance-test} closes to $7.1\times10^{-17}$ in an
independent high-precision evaluation (dps $=60$, root from the exact coefficient
polynomial, null vector re-derived). Driven through the repository's own
mode-vector routine the same identity closes only to $\approx1.9\times10^{-5}$.
The identity is scale-invariant in the null vector, so that residual measures the
accuracy of the null-vector \emph{direction} for a matrix of condition number
$\sim10^{22}$ --- it is the amplitude-normalisation defect U-3, not a defect of
the energy identity. Check D2 reports this honestly as \textsc{unresolved} rather
than loosening the tolerance until it passes.
\end{remark}

\section{Bulk positivity}
\label{sec:positivity}

\begin{proposition}[Positivity of the bulk dissipation]
\label{prop:positivity}
Under the stability conditions \eqref{eq:thermal-stability} of the
governing-equations note --- in particular $\kappa_{ij}\succ0$, $\rho>0$,
$c_E>0$, $T_0>0$ and $\tau_q\ge0$ --- the cycle average of
$\mathcal{D}_{\text{bulk}}$ in \eqref{eq:D-both} is non-negative.
\end{proposition}
\begin{proof}
Write the flux law as
$\bq = -\kappa\nabla\bth/(1-i\omega\tau_q) = -\Lam^{-1}\kappa\nabla\bth$ with
$\Lam = 1-i\omega\tau_q$. Then
$\bq\cdot\nabla\bth = -\Lam^{-1}(\nabla\bth)\cdot\kappa(\nabla\bth)$ and
$\tau_q\bq\cdot\partial_t\bq = \tau_q\bq\cdot(-i\omega\bq)
= i\omega\tau_q\lvert\bq\rvert^2_{\kappa^{-1}}$ (in the $\kappa$-weighted norm).
Substituting into \eqref{eq:D-both} and taking the real part of the cycle
average (with the $q_x$ term for the transverse direction),
\begin{equation}
  \langle\mathcal{D}_{\text{bulk}}\rangle
  = \frac{1}{T_0}\frac{\kappa}{1+\omega^2\tau_q^2}
    \big\lvert\nabla\bth\big\rvert^2 \;\ge\;0 ,
  \label{eq:bulk-diss-explicit}
\end{equation}
where the $1+\omega^2\tau_q^2$ factor comes from
$\operatorname{Re}(\Lam^{-1}) = 1/(1+\omega^2\tau_q^2)$ for
$\Lam=1-i\omega\tau_q$. Non-negativity follows from $\kappa>0$, $T_0>0$ and
$\tau_q\in\mathbb{R}$. Note that the relaxation term modifies the \emph{effective}
conductivity, $ \kappa/(1+\omega^2\tau_q^2)$, and does not add a second,
independent source of dissipation.
\end{proof}
\textsc{Derived}. The explicit form \eqref{eq:bulk-diss-explicit} is a useful
by-product: it shows that the Lord--Shulman relaxation \emph{reduces} the
effective conductivity at high frequency, which is the physically expected
behaviour and is a check on the sign convention $\Lam = 1-i\omega\tau_q$. Using
$1+i\omega\tau_q$ would give the same real part, so this particular check does not
distinguish the two conventions; the decay of the thermal branch does
(\texttt{characteristic\_matrix.tex}, the section on the roots).

\section{Interface dissipation}
\label{sec:interface}

The contact laws and their dissipations are derived in
\texttt{theory/interface\_conditions.tex}. For completeness the three forms are:
\begin{align}
  \mathcal{D}_{\mathrm{mech}} &= \tfrac12\,\jmp{\bv}^{\mathsf T}C\jmp{\bv}
     \;\ge\;0, \label{eq:D-mech}\\
  \mathcal{D}_{\mathrm{elec}} &= \tfrac12 G_s\lvert\jmp{\varphi}\rvert^{2}
     \;\ge\;0 \quad(\text{only if a leakage conductance }G_s\text{ is declared}),\\
  \mathcal{D}_{\mathrm{th}} &= \frac{H_T\lvert\jmp{\bth}\rvert^{2}}{T_0}
     \;\ge\;0, \label{eq:D-th}
\end{align}
all in $\mathrm{W\,m^{-2}}$. The capacitive gap of the companion note is
\emph{reactive} and contributes nothing to $\mathcal{D}_{\mathrm{elec}}$ unless
$G_s$ is declared.

\begin{remark}[Factor-of-two convention, still open]
\label{rem:half}
$\mathcal{D}_{\mathrm{mech}}$ and $\mathcal{D}_{\mathrm{th}}$ are written with a
$\tfrac12$ and without one respectively. \textbf{Phase 1 claimed this reflected a
genuine difference of origin; that was a misdiagnosis.} The difference is purely a
matter of whether the quantity is already a cycle average: both
$\tfrac12\operatorname{Re}(AB^*)$ forms carry the same $\tfrac12$, and
$\mathcal{D}_{\mathrm{th}}=H_T\lvert\jmp{\bth}\rvert^2/T_0$ is written for the
\emph{real} temperature jump field, not for a phasor. Written consistently in
phasors it becomes $\tfrac12 H_T\lvert\jmp{\Theta}\rvert^2/T_0$, which removes the
apparent asymmetry. The remaining question is only which convention the Phase-2
solver adopts for $R$, $T$, $A$, so U-5 is retained as a \emph{bookkeeping}
decision rather than a physics uncertainty. This is recorded as
U-5 in \texttt{reports/UNRESOLVED\_ISSUES.md} and is \textbf{not} resolved here,
because the correct choice depends on the normalisation adopted for $R$, $T$ and
$A$ in the Phase-2 solver.
\end{remark}

\section{The composite power balance}
\label{sec:balance}

Integrating \eqref{eq:local-balance} over the slab between two evaluation planes
$z=z_1<0$ and $z=z_2>\mathrm{h}$, and using the divergence theorem, gives
\begin{equation}
  \boxed{
  \underbrace{\langle J_z\rangle\big|_{z_1}}_{\text{incoming}}
  = \underbrace{\big(-\langle J_z\rangle\big|_{z_1}\big)_{\text{reflected}}}_{\text{outgoing at }z_1}
  + \underbrace{\langle J_z\rangle\big|_{z_2}}_{\text{transmitted}}
  + \underbrace{\int_{z_1}^{z_2}\langle\mathcal{D}_{\text{bulk}}\rangle\,dz}_{\text{bulk absorption}}
  + \underbrace{\sum_{\text{faces}}\langle\mathcal{D}_{\text{face}}\rangle}_{\text{interface absorption}}
  }
  \label{eq:master}
\end{equation}
\textsc{Derived}. Equation~\eqref{eq:master} is the form of Blueprint Eq.~(8) with
the terms identified explicitly. Dividing by the incoming flux defines
\begin{equation}
  R = \frac{\text{outgoing at }z_1}{\text{incoming}},\quad
  T = \frac{\text{transmitted at }z_2}{\text{incoming}},\quad
  A = \frac{\text{bulk + interface absorption}}{\text{incoming}},
  \label{eq:RTA}
\end{equation}
whence $R+T+A=1$ \emph{follows} from \eqref{eq:master}. It is not a premise.

\subsection{What must be fixed before $R,T,A$ may be quoted}

Blueprint \S13 item 4 requires the definition of reflected and transmitted power,
the evaluation planes, and the separation from bulk loss to be fixed
\emph{before} quoting $R$, $T$ or $A$. The list, in order:
\begin{enumerate}
  \item \textbf{Evaluation planes.} Both planes must lie in the homogeneous
        half-spaces, and $z_2$ must be far enough that the evanescent thermal and
        electrical modes have decayed. With $\mu\approx3.9\times10^{-7}$ m at
        1 MHz (\texttt{characteristic\_matrix.tex}), even $z_2=10\ \mu\mathrm{m}$
        is $25$ diffusion lengths --- ample. This must be verified, not assumed.
  \item \textbf{Channel-wise versus total.} At oblique incidence the reflected
        field contains two propagating channels (quasi-P and quasi-SV) whose
        fluxes are \emph{not} additive in $\lvert r\rvert^2$; the cross terms
        between them contribute. $R$ is defined from the \emph{total} reflected
        flux, not from a sum of $\lvert r_i\rvert^2$.
  \item \textbf{Bulk absorption inside the evaluation window.} If the outer media
        are lossy, $\int_{z_1}^{0}\langle\mathcal{D}_{\text{bulk}}\rangle dz$ is
        non-zero and belongs to $A$, not to $R$ or $T$. Leaving it out makes
        $R+T>1$ (or $<1$) and is the mechanism by which a pure-reflection
        optimiser can certify an attenuating window as ``anti-reflecting''.
  \item \textbf{The interphase.} For the primary configuration the interphase
        sits between the planes and its bulk absorption is part of $A$; the two
        face dissipations are separate contributions.
\end{enumerate}

\begin{remark}[A warning that follows from the derivation]
\label{rem:warning}
Because $A$ in \eqref{eq:RTA} collects both bulk and interface absorption,
a reduction in $R$ accompanied by an increase in $A$ is \emph{not} evidence of
anti-reflection with useful transmission. This is precisely the distinction
Blueprint \S2 rule F1 enforces, and the derivation here shows it is structural,
not a reporting preference: $R+T+A=1$ means any decrease in $R$ must appear
either in $T$ or in $A$, and separating the two requires the integrals in
\eqref{eq:master} to be evaluated, not inferred from $R$ alone.
\end{remark}

\section{What this note does and does not establish}
\label{sec:summary}

\paragraph{Established (\textsc{Derived}, and \textsc{Tested} where stated).}
The local balance \eqref{eq:local-balance} and the total energy current
\eqref{eq:total-current}; the \emph{corrected} cycle-averaged normal flux
\eqref{eq:Jz} with its bilinear thermal term
(Proposition~\ref{prop:flux}); its agreement with a direct cycle average of the
real fields; its exact antisymmetry; its $z$-conservation; the Second-Law sign of
the thermal term (Proposition~\ref{prop:thermal-sign}); the explicit bulk
dissipation \eqref{eq:bulk-diss-explicit} and its positivity
(Proposition~\ref{prop:positivity}); the interface forms
\eqref{eq:D-mech}--\eqref{eq:D-th}; and the composite balance
\eqref{eq:master}, from which $R+T+A=1$ follows.

\paragraph{Withdrawn in Phase 1B.}
The Phase-1 formula \eqref{eq:phase1-wrong} is \emph{withdrawn}. It had a linear
$+\tfrac12\operatorname{Re}[q_z]$ where a bilinear
$+\tfrac12\operatorname{Re}[q_z\Theta^*]/T_0$ is required, and it was wrong by up
to $34$ orders of magnitude on the thermal branch. Any downstream use of that
expression is invalid.

\paragraph{Not established.}
\begin{itemize}
  \item \textbf{No numerical energy balance.} Equations \eqref{eq:master} and
        \eqref{eq:RTA} are derived but have not been evaluated for any
        configuration, because the coupled mode-matching solver does not exist.
        Gate G5 (E3) is \textsc{unresolved}.
  \item \textbf{No partition of $A$ between bulk and interface} has been
        computed for any material pair.
  \item The factor-of-two convention of Remark~\ref{rem:half} is open (U-5).
  \item The claim that the thermal mode contributes negligible flux at practical
        evaluation planes is argued from the diffusion length
        $\mu\approx3.9\times10^{-7}$ m but not demonstrated numerically in a
        scattering configuration.
  \item Nothing in this note is a reflection-suppression result. No
        anti-reflection curve, and no claim about reflection minima, appears
        here.
\end{itemize}

\end{document}
